% Part: lambda-calculus
% Chapter: introduction
% Section: church-rosser

\documentclass[../../../include/open-logic-section]{subfiles}

\begin{document}

\olfileid{lam}{int}{cr}
\olsection{চার্চ--রসার ধর্ম}

\begin{thm}
\ollabel{thm:church-rosser}
ধরা যাক $M$, $N_1$ ও $N_2$ এমন পদ, যেখানে $M \red N_1$ এবং $M \red
N_2$। তাহলে এমন একটি পদ~$P$ আছে যাতে $N_1 \red P$ এবং $N_2 \red P$।
\end{thm}

\begin{cor}
ধরা যাক $M$-কে স্বাভাবিক রূপে হ্রাস করা যায়। তাহলে এই স্বাভাবিক রূপটি
একক।
\end{cor}

\begin{proof}
$M \red N_1$ এবং $M \red N_2$ হলে আগের উপপাদ্য অনুসারে এমন একটি
পদ~$P$ আছে যাতে $N_1$ ও $N_2$ উভয়ই~$P$-তে হ্রাস পায়। $N_1$ ও~$N_2$
উভয়ই স্বাভাবিক রূপে থাকলে এটি কেবল তখনই ঘটতে পারে যখন $N_1 \ident P
\ident N_2$।
\end{proof}

সবশেষে, দুটি পদ~$M$ ও~$N$ একটি সাধারণ পদে হ্রাস পেলে তাদের
\emph{$\beta$-সমতুল্য}, অথবা শুধু \emph{সমতুল্য}, বলব; অন্যভাবে বললে,
এমন কোনো~$P$ থাকতে হবে যাতে $M \red P$ এবং $N \red P$। এটি $M
\equal[\beta] N$ লিখি। \olref{thm:church-rosser} ব্যবহার করে যাচাই করা
যায় যে $\equal[\beta]$ একটি সমতুল্যতা-সম্বন্ধ; তার আরও ধর্ম হলো, প্রত্যেক
$M$ ও~$N$-এর জন্য $M \red N$ অথবা $N \red M$ হলে $M \equal[\beta]
N$। (আসলে দেখানো যায়, এই ধর্মবিশিষ্ট সমতুল্যতা-সম্বন্ধগুলির মধ্যে
$\equal[\beta]$-ই \emph{ক্ষুদ্রতম}।)

\end{document}
